\section{Future：Evaluation、Multimodal RAG 与 RAG 2.0}

最后部分把课程收束到系统研究议程：从零 joint pretraining 仍不足，scaling law、chunker、encoder、database、synthetic data 与 evaluation 都可学习或重新设计；RAG 还可跨越文本进入图像和多模态；最终目标是优化整个 cost--quality system。

\subsection{Open Questions：还有哪些层没有被优化}

Open Questions slide 询问 joint pretraining、scaling laws、reranker/chunker、bi-encoder 之外的表示、database convergence、synthetic data 与 evaluation。它说明“RAG 已解决”远非事实，现有系统仍大量依赖手工切分和弱代理指标。

\lecturefigure{slide-47-open-questions.jpg}{RAG 的 joint training、chunking、database 与 evaluation 开放问题}{Stanford 课堂视频 01:07:36}

\begin{knowledgebox}{RAG Evaluation 的分层矩阵}
\begin{tabularx}{\textwidth}{L{0.18\textwidth}>{\raggedright\arraybackslash}X >{\raggedright\arraybackslash}X}
\toprule
层级 & 指标示例 & 典型盲点 \\
\midrule
Retrieval & Recall@$k$, MRR, nDCG & relevance label 不等于 generator utility \\
Packing & evidence coverage、position & 忽略 duplicate 与 token budget \\
Generation & correctness、citation、faithfulness & 正确答案可能来自参数记忆 \\
System & latency、cost、freshness、access control & 平均值掩盖尾延迟与高风险错误 \\
\bottomrule
\end{tabularx}
\end{knowledgebox}

\teachervoice{讲者批评当前 RAG evaluation 很弱：downstream performance 看不到 retrieval 在哪里出错，把 retrieval accuracy 与 LM accuracy 做 harmonic mean 也缺少可靠数据集支持。}

\subsection{Vector Database 会不会并入普通 Database}

在口头开放问题中，讲者怀疑专用 vector database 是否会长期独立存在：BM25 更高效，reranker 可能补足语义；传统数据库正在加入向量检索能力，向量数据库厂商也在补充 sparse retrieval 与 metadata filtering。这是 2023 年的产业判断，应作为架构趋势讨论，而不能当作当前市场事实。

\begin{warningbox}{不要把课堂预测写成产品结论}
数据库实现、硬件、索引算法和供应商能力变化很快。讲义只保留抽象问题：未来系统是否需要统一管理 sparse、dense、metadata、transactions 与 access control，而不声称某类产品已消失或胜出。
\end{warningbox}

\subsection{Multimodal RAG}

Multimodal slide 展示 cross-modal retrieval augmentation 与 visual question answering pipeline。文本 query 可以检索图像、图像可以检索文本，或者 vision pipeline 先生成结构化描述，再把结果交给 frozen LM。

\lecturefigure{slide-48-multimodal-rag.jpg}{Cross-modal retrieval 与 multimodal retrieval-augmented LM}{Stanford 课堂视频 01:08:45}

\begin{knowledgebox}{读图：Multimodal RAG 需要对齐三种空间}
系统至少要对齐 query representation、retrieved modality representation 与 generator input representation。图像相似不等于回答相关，caption 也可能丢失视觉细节；因此需要 cross-modal retriever、modality adapter 和 task-level evaluation。
\end{knowledgebox}

\begin{warningbox}{视觉检索结果也可能成为错误证据}
相似图片可能来自不同时间、地点或对象，OCR/caption 可能出错。Multimodal grounding 需要 provenance、时间与实体核对，不能因为“有图”就提高证据等级。
\end{warningbox}

\teachervoice{Kiela 把 multimodality 视为自然延伸：为什么 RAG 只停留在 text？他引用 LENS 与 cross-modal work，强调 frozen LM 也能通过外部 vision pipeline 获得视觉能力。}

\subsection{RAG 2.0：Systems over Models}

最终 slide 用四点收束：systems over models、optimize it all、trade cost and quality、zero-shot domain generalization。RAG 2.0 不是一个固定图，而是一种 optimization stance：retriever、chunker、reranker、generator、index、data 与 serving 不应被永久分割。

\lecturefigure{slide-49-rag-2-0.jpg}{RAG 2.0：从 frozen components 转向 systems-over-models}{Stanford 课堂视频 01:10:12}

\begin{importantbox}{Cost--Quality 的系统目标}
可把部署目标抽象为
\[
\max_{\Theta,I,\pi}\;
Q(\Theta,I,\pi)
-\lambda_L L(\Theta,I,\pi)
-\lambda_C C(\Theta,I,\pi)
-\lambda_R R(\Theta,I,\pi),
\]
其中 $Q$ 是任务质量，$L$ 是 latency，$C$ 是 compute/storage cost，$R$ 是风险或不可信输出；$\pi$ 包含 retrieval policy。只优化 LM accuracy 不能决定系统最优解。
\end{importantbox}

\teachervoice{讲者把 RAG 2.0 概括为 ``systems over models''，甚至建议反向传播到 chunker。核心不是所有模块必须同时更新，而是任何手工边界都应接受“是否可学习、是否值得优化”的检验。}

\subsection{Closing Q\&A：Hybrid、Fine-tuning 与 Retrieval Hardware}

最后 Q\&A 给出三个实践判断。第一，BM25 可先 cast a wide net，再用 dense reranker 收窄。第二，domain adaptation 最终会组合 retrieval、fine-tuning 与 adapters，而不是永远二选一。第三，讲者提到 dedicated retrieval hardware 的产业可能性，但没有给出可公开验证的型号或 benchmark，因此只记录为 speaker speculation。

\begin{knowledgebox}{Retrieval 与 Fine-tuning 解决不同状态}
Fine-tuning 改变 parameters，适合行为、格式、领域表示与稳定技能；retrieval 改变 per-request evidence，适合频繁更新、私有数据与 attribution。RAG 2.0 可以在领域任务上 fine-tune 整个 retrieval-augmented system，同时保留可更新 index。
\end{knowledgebox}

\teachervoice{讲者明确说，最终大家可能会把 retrieval、end-to-end fine-tuning 和 adapters 全部结合起来。真正问题变成怎样高效组合，而不是争论只能选择一种。}

\subsection{Hallucination、Generic Error 与 Ground Truth}

课程最重要的 spoken-only 段落出现在最后四分钟。讲者不满领域把所有错误都叫 hallucination：普通计算错误、误解问题与 evidence contradiction 应分开；hallucination 至少需要相对于某个 ground truth 或 retrieved evidence 定义。

\begin{importantbox}{三类失败的区分}
\begin{tabularx}{\textwidth}{L{0.20\textwidth}>{\raggedright\arraybackslash}X >{\raggedright\arraybackslash}X}
\toprule
类型 & 判定 & 例子 \\
\midrule
Generic error & 输出不满足任务答案或计算规则 & 算错数、漏约束、格式错误 \\
Unsupported claim & 没有足够 evidence 支持 & index 无相关来源却给出确定陈述 \\
Evidence contradiction & 输出与选定 ground truth/source 冲突 & 文档说 A，答案声称非 A \\
\bottomrule
\end{tabularx}
\end{importantbox}

\begin{warningbox}{Ground truth 选择本身是系统决策}
Index 可以只允许 peer-reviewed papers，也可以包含 arXiv、企业文档或开放网页；不同 source policy 会给出不同“可信事实”。RAG 让这一选择显式化，但错误或恶意 source 仍会让 grounded answer 错误。
\end{warningbox}

\teachervoice{讲者说，frozen GPT-4 即使接入 RePlug 仍可能忽略 evidence；若把知识全部移到外部 index、模型只负责语言与 reasoning，理论上更容易强制 grounding，但这是方向性极限而非已经实现的保证。}

\teachervoice{他进一步说，creative writing 有时需要“hallucinate”，而 factual task 要更强 grounding；理想系统应有 grounding knob。Temperature 只改变 sampling distribution，即使低温也可能稳定地产生错误，不能充当 grounding control。}

\subsection{本章小结}

Future RAG 需要联合 evaluation、database、multimodal retrieval、chunking、instruction 与 serving。RAG 2.0 的最终观点是优化整套 system，并清楚指定 source of truth、grounding strength、cost 和 risk；检索与 fine-tuning 是互补工具，temperature 不是事实性开关。

